% ============================================================
\subsubsection{Selection and Filtering}
% ============================================================

Three access needs recur: selecting by label, selecting by position, and selecting by predicate. pandas keeps the first two on separate accessors, \texttt{loc} for labels and \texttt{iloc} for positions, and expresses the third through a boolean mask. Restricting rows is the pandas \texttt{WHERE}, and every later pattern builds on it.

\paragraph{Bracket indexing.} The bare \texttt{[]} operator on a DataFrame is overloaded by the type of its argument. A string selects a column and returns a Series, a list of strings selects several columns and returns a DataFrame, and a boolean Series filters rows. The overloading is convenient but ambiguous, and a key that could name either a row or a column is resolved as a column. For anything past a plain column pick, \texttt{loc} states the intent without ambiguity.

\begin{lstlisting}[style=python, caption={The overloaded bracket}]
df['temp']            # a Series (one column)
df[['city', 'temp']]  # a DataFrame (a column list)
df[df['temp'] > 97]   # row filter by boolean mask
\end{lstlisting}

\vspace{-1ex}

\begin{definition}[loc and iloc]
\texttt{loc} selects by label and \texttt{iloc} by integer position. Both take a row selector and an optional column selector, \texttt{df.loc[rows, cols]}, and each selector accepts a single key, a list, a slice, or a boolean mask. A \texttt{loc} label slice includes both endpoints, whereas an \texttt{iloc} position slice excludes the stop in the usual Python way.
\end{definition}

The inclusive label slice is a frequent trap. \texttt{df.loc['a':'c']} returns rows a through c inclusive, while \texttt{df.iloc[0:3]} returns three rows and stops before position 3.

\begin{lstlisting}[style=python, caption={Label access, position access, filtered selection}]
df.loc[3, 'temp']                          # row label 3, column 'temp'
df.loc[df['temp'] > 97, ['city', 'temp']]  # mask picks rows, list picks columns
df.iloc[0]                                 # first row by position
df.iloc[-1]                                # last row by position
\end{lstlisting}

\vspace{-2ex}

\paragraph{Boolean masks.} A comparison on a Series returns a boolean Series of the same length, and passing it to \texttt{loc} keeps the \texttt{True} rows. Compound conditions combine with the bitwise operators \texttt{\&}, \texttt{|}, and \texttt{\textasciitilde}, and each comparison must be wrapped in parentheses because the bitwise operators bind tighter than the comparisons. The Python keywords \texttt{and}, \texttt{or}, and \texttt{not} do not work on a Series, since they demand a single truth value and raise on an array.

\newpage

\begin{lstlisting}[style=python, caption={Combine masks with bitwise operators, parenthesize each side}]
mask = (df['temp'] > 97) & (df['city'] == 'Austin')
df.loc[mask]

df.loc[~df['city'].isin(['Dallas', 'Houston'])]   # NOT IN
\end{lstlisting}

\vspace{-3ex}

\paragraph{Membership and range.} \texttt{isin(values)} is a vectorized membership mask, the pandas \texttt{IN}, and negating it with \texttt{\textasciitilde} gives the \texttt{NOT IN}. \texttt{between(lo, hi)} tests a closed interval. Both read more cleanly than the equivalent chain of \texttt{|} clauses or the paired inequalities they stand in for.

\paragraph{The query string.} \texttt{query} filters from a string expression, naming columns bare and outer variables with a leading \texttt{@}. It reads well for a long compound predicate and spares the repetition of the frame name. Being parsed at call time, it runs slower than a mask and accepts only expressions it can interpret, so a mask stays the general tool.

\begin{lstlisting}[style=python, caption={query names columns bare and locals with @}]
threshold = 97
df.query('temp > @threshold and city == "Austin"')
\end{lstlisting}

\begin{keybox}[Chained assignment is unreliable]
A selection can return a view onto the original's data or an independent copy, and which one arises is not predictable from the expression. Assigning through a chained selection, \texttt{df[df['t'] > 100]['city'] = 'hot'}, may therefore write into a discarded temporary and leave the original untouched, which is what raises the \texttt{SettingWithCopyWarning}. The reliable form selects the rows and the column together in a single \texttt{loc}, so the write lands on the frame itself.
\end{keybox}

\vspace{1ex}

\begin{lstlisting}[style=python, caption={The chained write warns and may no-op, the single loc is definite}]
df[df['temp'] > 100]['city'] = 'hot'          # unreliable, warns
df.loc[df['temp'] > 100, 'city'] = 'hot'      # correct, one selection
\end{lstlisting}

\vspace{-2ex}
